% 論文節の日本語訳（読み物 PDF）用 preamble
% lualatex + jlreq
\documentclass[12pt,paper=a4]{jlreq}
\usepackage{luatexja-fontspec}
\usepackage{amsmath,amssymb,bm}
\usepackage{graphicx}
\usepackage{hyperref}
\usepackage{enumitem}
\usepackage{booktabs}
\usepackage{microtype}
\usepackage{xcolor}
\usepackage{caption}
\usepackage{tcolorbox}
\tcbuselibrary{breakable,skins}

\hypersetup{
  colorlinks=true,
  linkcolor=blue!60!black,
  urlcolor=blue!60!black,
  citecolor=blue!60!black,
}

\graphicspath{{images/}{./}}

\setlength{\parindent}{1em}
\setlength{\parskip}{0.55em}
\setlist[itemize]{leftmargin=1.6em,itemsep=0.25em,topsep=0.35em}
\setlist[enumerate]{leftmargin=1.8em,itemsep=0.25em,topsep=0.35em}

\captionsetup{
  font=small,
  labelfont=bf,
  skip=0.4em,
  justification=raggedright,
  singlelinecheck=false,
}

\newcommand{\vx}{\bm{x}}
\newcommand{\vu}{\bm{u}}
\newcommand{\veps}{\bm{\epsilon}}
\newcommand{\vw}{\bm{w}}

\DeclareMathOperator*{\argmax}{arg\,max}
\DeclareMathOperator*{\argmin}{arg\,min}
\DeclareMathOperator{\Tr}{Tr}
\DeclareMathOperator{\diag}{diag}

% 命題・定理ボックス（ページ跨ぎ時も切れ目を抑える）
\newtcolorbox{propositionbox}[1][]{%
  enhanced,
  breakable,
  colback=blue!4!white,
  colframe=blue!45!black,
  coltitle=white,
  fonttitle=\bfseries\small,
  title={#1},
  attach boxed title to top left={yshift=-2mm, xshift=3mm},
  boxed title style={
    colback=blue!45!black,
    colframe=blue!45!black,
    boxrule=0pt,
    arc=1pt,
  },
  boxrule=0.5pt,
  arc=1.5pt,
  left=5pt, right=5pt, top=8pt, bottom=5pt,
  before skip=0.9em,
  after skip=0.9em,
  pad at break=3mm,
}

% 図はフロートせず本文位置に置く（分断を防ぐ）
\newcommand{\translationfigure}[2]{%
  \begin{center}
    \includegraphics[width=0.88\linewidth]{#1}\\[0.35em]
    {\small #2\par}
  \end{center}
  \vspace{0.4em}%
}

\begin{document}
{\color{gray}\footnotesize 出典: §4.2\par}
\vspace{0.2em}

\section*{4.2 前向き／逆時間 SDE と PF-ODE における周辺の整合}


本節では、Fokker–Planck 方程式を拡散モデルを支配する中心原理として提示する。それは、ノイズのある動力学の下で確率密度がどのように発展するかを記述し、ランダム性が存在するときの変数変換公式の自然な類似物として機能する（第 B 章を参照）。周辺 $\{p_t\}$ を定義する前向き拡散過程から始め、逆時間 SDE と PF-ODE の双方が同じ周辺族を再現するよう設計される様子を説明する。次いで、逆時間 SDE と PF-ODE を直接かつ明示的に計算できるガウス例を通じて、この原理を具体化する。

\subsection*{4.2.1 周辺密度の整合を保証する Fokker–Planck 方程式}


拡散モデルにおける中心概念は、異なる過程が同じ周辺分布の系列をもたらしうることである（本小節の後半で図示する）。目標は、時間にわたって——とくに $t=0$ で——周辺を整合させることで $p_{\mathrm{prior}}$ を $p_{\mathrm{data}}$ へ変換する過程を構成することである。過程の厳密な形は二次的であり、扱いやすく効率的なサンプリングを認めることが条件である。これは自然に次の根本的な問いにつながる：

\begin{propositionbox}[問い 4.2.1]
異なる過程が同一の周辺分布を与えることを、どのように保証できるか？
\end{propositionbox}


設定に戻ると、いったん前向き SDE が指定されれば、それは $p_{\mathrm{data}}$ から $p_{\mathrm{prior}}$ への周辺密度の発展を定義する。逆時間 SDE と PF-ODE は、その軌跡が前向き過程のものと正確に一致する周辺分布を与えるよう構成される。この対応の鍵は、拡散過程の下で周辺密度がどのように発展するかを支配する Fokker–Planck 方程式にある。図示は図4.5に与える。次の定理（Anderson, 1982; Song et al., 2020c）がこの接続の基礎を確立する：

\translationfigure{fig-4-5.jpg}{図4.5: 同じ周辺 $\{p_t\}$ を共有する三つの動力学。上段：ノイズ付加方向。アンカー参照として機能する前向き SDE が、データをガウス事前へ拡散することで周辺 $\{p_t\}$ を定義する。PF-ODE は時間前向きに走らせると同じ周辺経路を決定論的に辿る。下段：デノイジング方向。逆時間 SDE と時間を遡って走る PF-ODE は、Fokker–Planck 方程式により保証されるように、逆向きに同じ周辺を与える。逆時間 SDE は確率的デノイジング過程を、PF-ODE は決定論的なものを与える。したがって繰り返される密度スナップショットは意図的である：過程は一時刻周辺のレベルでは一致するが、サンプルレベルの動力学では異なる。PF-ODE は決定論的流れでありいずれの時間方向にも積分できるため、両パネルに現れる。}


\begin{propositionbox}[定理 4.2.1: Fokker–Planck 方程式による周辺整合]
前向き SDE

\[
dx(t)=f\bigl(x(t),t\bigr)\,dt + g(t)\,dw(t), \qquad x(0)\sim p_0=p_{\mathrm{data}}
\]

を考える。標準的な正則性条件の下で、その周辺密度 $p_t$ は Fokker–Planck 方程式

\[
\begin{aligned}
\partial_t p_t(x)
&= -\nabla_x\cdot\bigl[f(x,t)p_t(x)\bigr] + \tfrac12 g^2(t)\Delta_x p_t(x) \\
&= -\nabla_x\cdot\bigl[v(x,t)p_t(x)\bigr] \quad (4.2.1)
\end{aligned}
\]

に従って発展する。ここで $\Delta_x$ はラプラシアン作用素であり、

\[
v(x,t):=f(x,t)-\tfrac12 g^2(t)\nabla_x\log p_t(x).
\]

重要なことに、同じ周辺分布の族 $\{p_t\}_{t\in[0,T]}$ は二つの異なる動力学によって実現できる：

\textbf{(i)} PF-ODE

\[
\frac{d\tilde{x}(t)}{dt}=v\bigl(\tilde{x}(t),t\bigr).
\]

$\tilde{x}(0)\sim p_0$ から始め、$0$ から $T$ へ前向きに解くと $\tilde{x}(t)\sim p_t$ が得られる。等価に、$\tilde{x}(T)\sim p_T$ から始め同じ ODE を $T$ から $0$ へ後ろ向きに解くと、同じ周辺を逆に辿る。

\textbf{(ii)} 逆時間 SDE

\[
d\bar{x}(t)=\Bigl[f\bigl(\bar{x}(t),t\bigr)-g^2(t)\nabla_x\log p_t\bigl(\bar{x}(t)\bigr)\Bigr]dt + g(t)\,d\bar{w}(t).
\]

$\bar{x}(T)\sim p_T$ から始め、この SDE は時間を $t=T$ から $t=0$ へ遡って解かれる。この逆発展に沿って、時刻 $t$ におけるその周辺分布は

\[
\bar{x}(t)\sim p_t
\]

である。

したがって、逆時間 SDE と PF-ODE は、一方が確率的で他方が決定論的であるにもかかわらず、すべてのノイズレベルで同じ周辺分布をもちうる。
\end{propositionbox}


\textbf{定理の証明。}

証明は第 D.2.5 節に与えられ、第4.6.1節は確率の周辺化技法を用いて Fokker–Planck 方程式の背後のさらなる直観を提供する。■

\textbf{固定周辺に対する複数の条件付き分布。} PF-ODE が $p_{\mathrm{data}}$ を時間前向きに（あるいは等価に $p_{\mathrm{prior}}$ を逆に）輸送する様子を理解するため、流れ写像 $\Psi_{s\to t}:\mathbb{R}^D\to\mathbb{R}^D$ を考える。ここで $\Psi_{s\to t}(x_s)$ は、任意の時刻 $s,t\in[0,T]$ について、時刻 $s$ における $x_s$ から初期化された PF-ODE 解の時刻 $t$ における値を表す。言い換えれば、この写像は初期状態 $x_s$ をとり、その時刻 $t$ における状態へ直接ジャンプする：

\[
\Psi_{s\to t}(x_s) := x_s + \int_s^t v(x_\tau,\tau)\, d\tau, \quad (4.2.2)
\]


速度場は

\[
v(x,\tau):=f(x,\tau)-\tfrac12 g^2(\tau)\nabla_x\log p_\tau(x).
\]


ここで積分は PF-ODE 軌跡に沿って蓄積された正味変位を捉える。流れ写像という語は実際、微分方程式、力学系、連続体力学において古典的である：それは与えられた動力学の下で、状態を時刻 $s$ における位置から時刻 $t$ における位置へ運ぶ写像を表す。速度場 $v(x,t)$ は瞬間運動を記述し、流れ写像 $\Psi_{s\to t}$ は有限時間区間にわたる蓄積運動を記録する（arnold1988geometrical; chorin1993mathematical）\footnote{この写像の根底にある物理的考えははるかに古い。18世紀の流体力学はすでに、速度場を通じ、また物質粒子の発展を追うことで運動を研究していた。とくに Euler と Lagrange の基礎的著作（euler1757fluid; lagrange1783fluid）においてである。}。本書を通じて、とくに第10章および第11章で高速な拡散ベース生成器を議論するとき、この標準的な数学的・物理的意味でこの語を用いる。$v$ に関する穏やかな滑らかさ仮定の下で、流れ写像 $\Psi_{s\to t}:\mathbb{R}^D\to\mathbb{R}^D$ は滑らかな全単射である\footnote{ネタバレ：PF-ODE 流れ写像 $\Psi_{s\to t}$ は、$p_s$ を $p_t$ へ運ぶ正規化フロー（Normalizing Flow; NF）の全単射そのものである（第5.1節で詳述）。違いは、いったん前向き SDE とその周辺密度が固定されれば、PF-ODE 公式がその関連速度場 $v(x,t)=f(x,t)-\tfrac12 g^2(t)\nabla_x\log p_t(x)$ を一意に指定するのに対し、正規化フロー、あるいは連続時間正規化フローは速度場を直接パラメータ化して学習することである。双方とも同じ連続変数変換原理に依拠する。}。

任意の $t\in[0,T]$ について、押し出し密度は

\[
p_t^{\mathrm{fwd}}(x_t) := \int \delta\bigl(x_t-\Psi_{0\to t}(x_0)\bigr)\, p_{\mathrm{data}}(x_0)\, dx_0
\]


と定義され、$\Psi_{0\to t}\sharp p_{\mathrm{data}}$ と記され、$\Psi_{0\to t}$ の下での時刻 $t$ における分布を表す。定理 4.2.1 は $p_t^{\mathrm{fwd}}=p_t$ を保証する。ここで $p_t$ は前向き SDE の周辺密度であり、決定論的 PF-ODE と確率的核を等置する：

\[
p_t(x_t)=\int p_t(x_t\mid x_0)\, p_{\mathrm{data}}(x_0)\, dx_0
=\int \delta\bigl(x_t-\Psi_{0\to t}(x_0)\bigr)\, p_{\mathrm{data}}(x_0)\, dx_0.
\]


これは、無限に多くの条件付き $Q_t(x_t\mid x_0)$ が同じ $p_t(x_t)$ を与えることを含意する。例えば：

\begin{itemize}
\item \textbf{確率的（シミュレーションフリー）}: $Q_t(x_t\mid x_0)=p_t(x_t\mid x_0)$,
\item \textbf{決定論的（ODE 求解が必要）}: $Q_t(x_t\mid x_0)=\delta\bigl(x_t-\Psi_{0\to t}(x_0)\bigr)$,
\item \textbf{混合}: $Q_t(x_t\mid x_0)=\lambda p_t(x_t\mid x_0)+(1-\lambda)\delta\bigl(x_t-\Psi_{0\to t}(x_0)\bigr)$, $\lambda\in[0,1]$.
\end{itemize}


この $Q_t(x_t\mid x_0)$ の非一意性は、周辺制約が条件付き分布を一意に定めないという事実から生じる。この概念は第5.2.2節および第9.2.3節に再び現れる。とくに、同じ周辺 $p_t$ と整合する逆時間 SDE の族全体が存在する。

\begin{propositionbox}[観察 4.2.1: 規定された周辺密度の整合]
複数の過程が同じ周辺密度の系列を生じうる；真に重要なのは Fokker–Planck 方程式を満たすことである。この根本的洞察は、$p_{\mathrm{prior}}$ から $p_{\mathrm{data}}$ へ、あるいはその逆へ遷移する生成過程の設計において、顕著な柔軟性を与える。
\end{propositionbox}


Fokker–Planck 方程式は拡散モデルの中心にあり、確率密度の基本的な変数変換公式に根ざす（系統的取り扱いについては第 B 章を参照）。些細な技術的詳細などではなく、この原理は我々の展開を通じて繰り返し現れ、とくに第5.2節で顕著である。


\subsection*{4.2.2 計算可能な例：ガウス動力学の発展}


$p_{\mathrm{data}}$ が正規分布（またはガウス混合）であるとき、スコア関数は閉形式表現を認める。これは拡散過程についての直観を築く理想的な設定である：高度な数学的道具に頼ることなく、基本的な微積分のみを用いて逆時間 SDE と PF-ODE を明示的に導出できる。本小節では、そのような扱いやすい場合にこれらの方程式がどのように振る舞うかを示す。

\textbf{ガウスによる逆時間 SDE の厳密計算。} $p_{\mathrm{data}}$ がガウスであるとき、式 (4.1.6) の公式は、Anderson (1982) の一般理論と証明に依拠することなく直接導出できる。核心の考えを示すため、一次元の場合を考える；高次元への拡張は同様に従う。

前向き SDE

\[
dx(t)=f(t)x(t)\,dt + g(t)\,dw_t
\]


から始め、サイズ $\Delta t>0$ の小さな Euler ステップを一つとる：

\[
x_{t+\Delta t}=a x_t + r\epsilon,
\]


ここで $a:=1+f(t)\Delta t$、$r:=g(t)\sqrt{\Delta t}$、$\epsilon\sim\mathcal{N}(0,1)$ である。等価に、前向き一段遷移核はガウスである：

\[
x_{t+\Delta t}\mid x_t \sim \mathcal{N}(\,\cdot\,;\, a x_t,\, r^2).
\]


$p_{\mathrm{data}}$ がガウスであると仮定されているため、時刻 $t$ における現在の周辺もガウスであり、次の形をとる：

\[
x_t\sim\mathcal{N}(\,\cdot\,;\, m_t,\, s_t^2),
\]


あるスカラー $m_t$ と $s_t$ について。したがって条件付けは二つのガウスの積と再正規化に帰着し、代数を初等的に保つ。

ベイズの法則により、条件付き密度は定数を除き事前と遷移核の積である：

\begin{align*}
\begin{aligned}
p(x_t\mid x_{t+\Delta t})
&\propto p(x_{t+\Delta t}\mid x_t)\, p_t(x_t) \\
&\propto \exp\Bigl(-\frac{(x_t-m_t)^2}{2s_t^2}\Bigr)
\exp\Bigl(-\frac{(x_{t+\Delta t}-a x_t)^2}{2r^2}\Bigr).
\end{aligned}
\end{align*}


指数は $x_t$ の二次式である。両平方を展開し項をまとめると、どの係数が重要かが正確に示される：

\[
-2\log p(x_t\mid x_{t+\Delta t}) = A x_t^2 - 2B x_t + \mathrm{const},
\]


ここで

\[
A:=\frac{1}{s_t^2}+\frac{a^2}{r^2}, \qquad
B:=\frac{m_t}{s_t^2}+\frac{a\, x_{t+\Delta t}}{r^2}.
\]


ここで $A$ は精度の和（事前精度に加え、$a$ を通じて輸送された遷移核精度）であり、$B$ は対応する精度重み付きターゲットの和である。これらが手元にあれば、平方完成が一行で事後を与える：

\[
A x_t^2 - 2B x_t = A\Bigl(x_t-\frac{B}{A}\Bigr)^2 - \frac{B^2}{A},
\]


したがって条件付き分布は分散 $1/A$、平均 $B/A$ のガウスである：

\[
\mathrm{Var}(x_t\mid x_{t+\Delta t})=\Bigl(\frac{1}{s_t^2}+\frac{a^2}{r^2}\Bigr)^{-1}, \qquad
\mathbb{E}[x_t\mid x_{t+\Delta t}]=\frac{\frac{m_t}{s_t^2}+\frac{a\, x_{t+\Delta t}}{r^2}}{\frac{1}{s_t^2}+\frac{a^2}{r^2}}.
\]


これらの閉形式は、任意の小さな $\Delta t$ に対する逆遷移をすでに記述する。逆時間 SDE を読み取るため、小さな $\Delta t$ についてそれらを展開する。

$a=1+f(t)\Delta t$ および $r^2=g^2(t)\Delta t$ を用いる。$\Delta t\to 0$ につれ、寄与 $a^2/r^2\sim 1/(g^2(t)\Delta t)$ が精度を支配するので、分散は

\[
\mathrm{Var}(x_t\mid x_{t+\Delta t})=\Bigl(\frac{1}{s_t^2}+\frac{a^2}{r^2}\Bigr)^{-1}=g^2(t)\Delta t+O(\Delta t^2)
\]


となり、逆ステップが前向きステップと同じ拡散スケール $g(t)$ をもつことを告げる。平均については、比 $B/A$ を一次まで展開する：

\[
\mathbb{E}[x_t\mid x_{t+\Delta t}]
= x_{t+\Delta t}
+ \Delta t\Biggl[
-\Bigl(f(t)+\frac{g^2(t)}{s_t^2}\Bigr)x_{t+\Delta t}
+ \frac{g^2(t)}{s_t^2}m_t
\Biggr]
+ O(\Delta t^2).
\]


平均と分散を合わせると、一段逆遷移核が得られる：

\[
x_t\mid x_{t+\Delta t}
\sim \mathcal{N}\Biggl(
x_{t+\Delta t}
+ \Delta t\Biggl[
-\Bigl(f+\frac{g^2}{s_t^2}\Bigr)x_{t+\Delta t}
+ \frac{g^2}{s_t^2}m_t
\Biggr],\,
g^2\Delta t+O(\Delta t^2)
\Biggr).
\]


これは、$t+\Delta t$ から $t$ へ後ろ向きに走る Euler–Maruyama 更新として認識される：

\[
x_t - x_{t+\Delta t}
= \Delta t\Biggl[
-\Bigl(f+\frac{g^2}{s_t^2}\Bigr)x_{t+\Delta t}
+ \frac{g^2}{s_t^2}m_t
\Biggr]
+ g\sqrt{\Delta t}\,\epsilon.
\]


この後ろ向きドリフトは、ガウス周辺 $p_t=\mathcal{N}(m_t,s_t^2)$ について

\[
\partial_x\log p_t(x)=-\frac{x-m_t}{s_t^2}
\implies
-\Bigl(f+\frac{g^2}{s_t^2}\Bigr)x+\frac{g^2}{s_t^2}m_t
= -\bigl[f x - g^2\partial_x\log p_t(x)\bigr]
\]


であるため、スコアで書ける。したがって $\Delta t\to 0$ とすると、元の拡散時計上の逆時間 SDE

\[
d\bar{x}(t)=\Bigl[f(t)\bar{x}(t)-g^2(t)\partial_x\log p_t\bigl(\bar{x}(t)\bigr)\Bigr]dt + g(t)\,d\bar{w}(t),
\qquad t:T\to 0
\]


が得られる。これは式 (4.1.6) の一般逆時間 SDE と正確に一致する。ベクトル形では、スカラー導関数 $\partial_x\log p_t$ は $\nabla_x\log p_t$ で置き換えられる。

\textbf{ガウスによる PF-ODE の厳密計算。} データ分布がガウスであると仮定するとき、Fokker–Planck 方程式のような重い道具を避けて、PF-ODE 公式も直接導出できる。最終的に、PF-ODE の周辺密度が前向き SDE および逆時間 SDE の双方のものと一致することを見、第4.2節で議論される Fokker–Planck 理論の構成的検証を与える。

時刻 $t$ で $x_t\sim\mathcal{N}(\,\cdot\,;\, m_t,\, s_t^2)$ と仮定する。サイズ $\Delta t$ の小さな決定論的ステップは、滑らかな写像

\[
x_{t+\Delta t}=\Phi_{t,\Delta t}(x_t)=x_t+\Delta t\, v_t(x_t)+O(\Delta t^2)
\]


として書け、これは単に $\Delta t$ に関する一次 Taylor 展開である。目標は、その決定論的流れが同じガウス周辺発展を再現する速度場を構成することである。とくに単純な選択はアフィン速度場

\[
v_t(x)=a_t x + b_t
\]


である。

アフィン動力学はガウス性を保存するため、平均と分散の発展を前向き SDE のものと整合させることで $a_t$ と $b_t$ を定められる。これをステップに戻すと

\[
x_{t+\Delta t}=(1+\alpha_t\Delta t)x_t+\beta_t\Delta t+O(\Delta t^2),
\qquad \alpha_t:=a_t,\ \beta_t:=b_t
\]


が得られる。

今、$x_t\sim\mathcal{N}(\,\cdot\,;\, m_t,\, s_t^2)$ をこの写像に通し、平均と分散を一次まで追跡する：

\begin{align*}
\begin{aligned}
\mathbb{E}[x_{t+\Delta t}] &= m_t + \Delta t(\alpha_t m_t + \beta_t) + O(\Delta t^2), \\
\mathrm{Var}(x_{t+\Delta t}) &= s_t^2 + \Delta t(2\alpha_t s_t^2) + O(\Delta t^2).
\end{aligned}
\end{align*}


他方、前向き SDE $dx=f(t)x\,dt+g(t)\,dw_t$ は初等的なモーメント公式をもつ（式 (4.4.3) を参照）：

\[
m_t'=f(t)m_t, \qquad (s_t^2)'=2f(t)s_t^2+g^2(t).
\]


$\Delta t$ の係数を整合させると

\[
\alpha_t=f(t)+\frac{g^2(t)}{2s_t^2}, \qquad
\beta_t=-\frac{g^2(t)}{2s_t^2}m_t
\]


が得られる。これらの選択で、ステップは

\[
x_{t+\Delta t}
= x_t
+ \Delta t\Biggl[
\Bigl(f(t)+\frac{g^2(t)}{2s_t^2}\Bigr)x_t
- \frac{g^2(t)}{2s_t^2}m_t
\Biggr]
+ O(\Delta t^2)
\]


となる。

ガウス $p_t=\mathcal{N}(m_t,s_t^2)$ について $\partial_x\log p_t(x)=-(x-m_t)/s_t^2$ であるため、括弧は $f(t)x_t-\tfrac12 g^2(t)\partial_x\log p_t(x_t)$ と書き換えられる。したがって、

\[
x_{t+\Delta t}
= x_t
+ \Delta t\Bigl[f(t)x_t-\tfrac12 g^2(t)\partial_x\log p_t(x_t)\Bigr]
+ O(\Delta t^2).
\]


最後に $\Delta t$ で割り $\Delta t\to 0$ とすると、PF-ODE

\[
x'(t)=f(t)x(t)-\tfrac12 g^2(t)\partial_x\log p_t\bigl(x(t)\bigr)
\]


が得られる。

なぜこの ODE が前向き SDE（および逆時間 SDE）と同じ周辺をもつかを見るには、上のドリフトが線形にシフトを加えたものであることに注意せよ。したがって $x(t)$ は $x(0)$ にアフィンに依存し、アフィン写像はガウスをガウスへ送る。さらに、この ODE に沿った平均 $m_t$ と分散 $s_t^2$ は、前向き SDE と正確に同じ二つのスカラー ODE を、同じ初期値で満たす（我々の整合による）。ゆえに $p_t=\mathcal{N}(m_t,s_t^2)$ は両発展についてすべての時刻 $t$ で同一である。
\end{document}
